\section{D'où vient l'argent}
\label{sec:carry}

Sur la sleeve or, cette section s'intitulait «~où passe l'argent~»~: le portage
y est un coût, et il fallait montrer combien il prélève. Ici la question
s'inverse. Le différentiel de taux entre le dollar et le yen a été positif sur
toute la fenêtre, et un long USD/JPY l'encaisse chaque nuit. Le poste change de
signe --- et de rang.

\begin{figure}[H]
\centering
\includegraphics[width=\linewidth]{usdjpy_trades_attribution}
\caption{Du résultat de prix au résultat net. Les trois postes sont mesurés sur
le journal des transactions et se recollent exactement au chiffre du
simulateur.}
\end{figure}

% Généré par scripts/build_usdjpy_trades_figures.py — NE PAS ÉDITER À LA MAIN.
\begin{table}[H]
\centering
\begin{tabular}{@{}lrl@{}}
\toprule
Poste & USD & Part \\
\midrule
Résultat de prix (mesuré) & +64,557 & 36.5\,\% du net \\
Portage (mesuré) & +112,266 & 63.5\,\% du net \\
Commission (mesurée) & +0 & nulle sur ce compte \\
\textbf{Résultat net} & \textbf{+176,822} & recolle au rapport MT5 \\
\bottomrule
\addlinespace[0.4em]
\multicolumn{3}{@{}p{0.92\linewidth}@{}}{\footnotesize Le portage vaut 0.69\,pb par jour de notionnel porté et retourne le signe de 3 trades, perdants sur le prix et gagnants une fois le carry compté. Slippage estimé à 25,350\,\$ (2\,bps par côté sur 63.4\,M\$ de notionnel cumulé) : il est déjà contenu dans les prix d'exécution, ce n'est pas un poste additionnel. La dernière position a été liquidée par le testeur en fin de fenêtre pour +37,343\,\$ : ce n'est pas une sortie décidée par le signal.}
\end{tabular}
\caption{Décomposition du résultat : ce qui est mesuré, ce qui est estimé}
\end{table}


Le résultat de prix --- ce que le mouvement du change a rapporté --- s'élève à
$+$64\,557\,\$. Le portage y ajoute $+$112\,266\,\$, soit \metric{174\,\%} de
plus, et la commission est nulle sur ce compte. Il reste $+$176\,822\,\$, qui
est exactement le chiffre du rapport d'exécution~: l'attribution est complète,
sans résidu.

\begin{insightbox}
\textbf{Ce candidat n'est pas seulement un momentum~: c'est un momentum
\emph{plus} un portage.} Le portage représente \metric{63\,\%} du résultat net.
Il retourne même le signe de \textbf{trois} trades, perdants sur le prix et
gagnants une fois le carry compté. Un instrument dont le classement repose à
63\,\% sur un poste que le signal ne pilote pas ne se juge pas comme les
autres.
\end{insightbox}

\subsection{Le portage est un loyer de temps --- perçu, pas payé}

\begin{figure}[H]
\centering
\includegraphics[width=\linewidth]{usdjpy_trades_costs}
\caption{Résultat de prix et résultat net cumulés~: l'écart entre les deux
courbes est le portage encaissé. À droite, le portage perçu par trade en
fonction de sa durée, sur deux échelles logarithmiques.}
\end{figure}

Le panneau de droite est presque une droite~: le portage croît proportionnellement
à la durée de détention. Rapporté au notionnel effectivement porté, il vaut
\metric{0,69~point de base par jour}, soit environ 2,5\,\% l'an du notionnel.
Multiplié par le levier moyen du moteur (6,1$\times$) et par son taux
d'exposition (67\,\%), cela représente de l'ordre de \metric{10\,\% l'an} du
sous-compte --- perçu sans qu'aucune décision de trading n'intervienne. Le
calcul direct donne la même chose~: 24\,276\,\$ de portage par an pour une
balance moyenne de 249\,289\,\$.

C'est ce qui explique la forme du panneau de gauche~: la courbe nette s'écarte
vers le haut de la courbe de prix et ne la rejoint jamais. Du sommet de
juillet~2024 au creux de février~2026, le résultat de prix cumulé perd
117\,704\,\$~; le résultat net n'en perd que 82\,355. Le portage a absorbé
\metric{30\,\%} des dégâts du signal sur cette phase.

\subsection{Trois réserves sur ce poste}

\textbf{Le carry est un régime, pas une constante.} La fenêtre 2021--2026
couvre l'épisode où la Réserve fédérale est montée au-delà de 5\,\% pendant que
la Banque du Japon restait à zéro. C'est le différentiel de taux le plus large
depuis des décennies, et rien dans ces trente-cinq trades ne dit ce que devient
le candidat si l'écart se referme. La question n'est pas hypothétique~: elle
porte sur les deux tiers du résultat.

\textbf{Le glissement pèse ici bien plus lourd que sur l'or.} Il est estimé à
\metric{25\,350\,\$} sur 63,4~millions de dollars de notionnel cumulé, à partir
de l'hypothèse de 2~points de base par côté --- soit \metric{39\,\%} du
résultat de prix. Cette somme est déjà contenue dans les prix d'exécution du
tableau ci-dessus~: ce n'est pas un poste à retrancher en plus. Mais si
l'hypothèse était deux fois trop optimiste, le résultat de prix perdrait encore
25\,000\,\$ et le résultat net 14\,\%. Sur l'or, la même hypothèse ne
représentait que 1,2\,\% du résultat de prix~; le candidat USD/JPY est
nettement plus sensible à la qualité d'exécution supposée.

\textbf{Le résultat de prix ne suffirait pas.} Sans portage, ce candidat aurait
rapporté $+$64\,557\,\$ au lieu de $+$176\,822\,\$, pour un risque au moins
aussi élevé --- le repli d'équité de 45\,\% est mesuré \emph{portage compris}.
Il est donc honnête de dire que c'est le \emph{carry} qui a placé USD/JPY en
tête du cycle, autant que le signal.
